\documentclass{ximera}
\title{Exterior derivatives}

\begin{document}

    \begin{abstract} Exterior derivatives, charts,  Leibniz rule, and nilpotency.  \end{abstract}
    \maketitle

        A reference for this material is Chapter 14 of John M. Lee. Introduction to smooth manifolds. Second edition. Grad. Texts in Math., 218. Springer, New York, 2013. xvi+708pp. ISBN: 978-1-4419-9981-8.

    \begin{definition}[Exterior derivative]
        For $\omega \in \Omega ^k (M)$, its \emph{exterior derivative} $d\omega \in \Omega ^{k + 1} (M)$ is given by
        \begin{align*}
            d\omega (X_0, \ldots , X_k) : = & \sum _{i = 0} ^k  (-1) ^i   X_i \omega (X_0, \ldots , \hat{X}_i , \ldots , X_k) \\
            & + \sum_{0\leq i < j \leq k} (-1) ^{i + j} \omega ( [X_i, X_j], X_0, \ldots , \hat{X}_i, \ldots , \hat{X}_j , \ldots , X_k  ),
        \end{align*}
        where the hat represents that the entry is missing. 
    \end{definition}

\begin{remark}
    If $f \in \Omega ^0 (M)$ is just a function, then 
    \[    d f (X) = X (f) .          \]
    In local coordinates, 
    \[     d f = \sum_{i = 1} ^n \frac{\partial f}{\partial x^i} dx ^i .       \]
\end{remark}
    


    \begin{proposition}[Well defined?] 
        $d\omega$ is indeed a  $(k+1 )$-form. 
    \end{proposition}

    \begin{solution}
        Exercise. Do it at least for $k =1$!
    \end{solution}


    \begin{proposition}[Exterior derivative in charts]
        If 
        \[   \omega = f \, dx^{i_1} \wedge \cdots \wedge dx^{i_k}    ,        \]
        then
        \[             d\omega =  df \wedge  dx^{i_1} \wedge \cdots \wedge dx^{i_k}   .                \]
        By the above remark, another way to write this is 
        \[   d\omega =  \sum _{j = 1 }^n \,\frac{\partial f }{\partial x ^j }\, dx^j \wedge  dx^{i_1} \wedge \cdots \wedge dx^{i_k}   .           \]
    \end{proposition}

    \begin{solution}
        Since the vector fields $ \tfrac{\partial }{\partial x ^1}, \ldots , \tfrac{\partial }{\partial x ^n} $ commute (partial derivatives commute), when we plug them into an exterior derivative, only the terms in the first sum show up.  For $1 \leq j_0 < \ldots < j_k \leq n$, we have

    \begin{align*}
        d \omega  (\tfrac{\partial }{\partial x ^{j_0}},  & \ldots ,  \tfrac{\partial }{\partial x ^{j_k}} )       =  \sum_{\ell = 0} ^k (-1 ) ^{\ell} \frac{\partial }{\partial x ^{j_{\ell}}}  \omega (\tfrac{\partial }{\partial x ^{j_0}}, \ldots ,\hat{\tfrac{\partial}{\partial x ^{j_{\ell}}}} , \ldots , \tfrac{\partial }{\partial x ^{j_k}} )   \\
        & = \sum _{\ell  = 0} ^k (-1 ) ^{\ell} \frac{\partial f }{\partial x ^{j_{\ell}}} dx ^{i_1} \wedge \cdots \wedge dx ^{i_k} ( \tfrac{\partial}{\partial x ^{j_0}} , \ldots , \hat{\tfrac{\partial}{\partial x ^{j_{\ell}}}} , \ldots ,  \tfrac{\partial}{\partial x ^{j_k}} ) 
     \end{align*}
    This sum is zero unless 
     \[    \{ i_1, \ldots , i_k \} \subset \{ j_0, \ldots , j_k \}   .                        \] 
    In that case, the only summand that is not zero is the one with 
    \[     \{ j _{\ell} \}  = \{ j_0, \ldots    , j_k \}  \backslash  \{ i_1, \ldots , i_k \}   ,  \] 
     and that summand is 
     \[    (-1) ^{\ell} \frac{\partial f}{\partial x ^{j_{\ell}}} .                      \]
     Plugging $(\tfrac{\partial }{\partial x ^{j_0}}, \ldots ,  \tfrac{\partial }{\partial x ^{j_k}} )$ into the right hand side of the expression of the proposition, we get zero unless 
     \[    \{ i_1, \ldots , i_k \} \subset \{ j_0, \ldots , j_k \} .          \]
     In that case, if we set
     \[ \{j_{\ell} \} =   \{ j_0, \ldots , j_k \} \backslash  \{ i_1, \ldots , i_k \} , \] 
     the only summand that is not zero is 
     \[      \frac{\partial f}{\partial x ^{j_{\ell}}}  d x ^{j_{\ell}} \wedge d x ^{i_1 } \wedge \cdots \wedge dx ^{i_k} ( \tfrac{\partial }{\partial x ^{j_0}}, \ldots ,  \tfrac{\partial }{\partial x ^{j_k}}  ) = (-1 )^{\ell} \frac{\partial f}{\partial x ^{j_{\ell}}} .      \]
     In either case, the two expressions agree.
    \end{solution}


    \begin{proposition}[Leibniz rule]
        For $\omega \in \Omega ^k (M)$ and $\eta \in \Omega ^{\ell} (M)$, one has 
        \[    d (\omega \wedge \eta  ) = (d \omega ) \wedge \eta + (-1 ) ^k \omega \wedge  (d \eta).          \]
    \end{proposition}
    \begin{solution} By linearity, it is enough to check in coordinates. Namely, 
    \begin{align*}
        \omega & = f \, dx^{i_1} \wedge \cdots \wedge dx ^{i_k}\\
        \eta & = g \, dx^{j_1} \wedge \cdots \wedge dx ^{j_{\ell}}.
    \end{align*}
    In that case, 
    \[      \omega \wedge \eta = fg \, dx^{i_1} \wedge \cdots \wedge dx ^{i_k} \wedge dx^{j_1} \wedge \cdots \wedge dx ^{j_{\ell}}.                     \]
    Taking the derivative, 
    \begin{align*}
        d (\omega \wedge \eta ) & = \sum _{a = 1 }^n \frac{\partial (fg)}{\partial x ^a} dx ^a \wedge dx^{i_1} \wedge \cdots \wedge dx ^{i_k} \wedge dx^{j_1} \wedge \cdots \wedge dx ^{j_{\ell}} \\
        & = \sum _{a = 1 }^n   \frac{\partial f }{\partial x ^a} g \,   dx ^a \wedge dx^{i_1} \wedge \cdots \wedge dx ^{i_k} \wedge dx^{j_1} \wedge \cdots \wedge dx ^{j_{\ell}} \\
        &  \,\,\,\,\,\,\,\, +  \sum _{a = 1 }^n  f \frac{\partial g}{\partial x ^a}  \,  dx ^a \wedge dx^{i_1} \wedge \cdots \wedge dx ^{i_k} \wedge dx^{j_1} \wedge \cdots \wedge dx ^{j_{\ell}} \\
        & = \left[ \sum_{a = 1}^n  \frac{\partial f }{\partial x ^a}    dx ^a \wedge dx^{i_1} \wedge \cdots \wedge dx ^{i_k}  \right] \wedge \left[ g   dx^{j_1} \wedge \cdots \wedge dx ^{j_{\ell}}  \right]\\
        & \,\,\,\,\,\, + (-1) ^k  \left[ f  dx^{i_1} \wedge \cdots \wedge dx ^{i_k}   \right] \wedge \left[ \sum _{a = 1 }^n   \frac{\partial g}{\partial x ^a}  \,  dx ^a  \wedge dx^{j_1} \wedge \cdots \wedge dx ^{j_{\ell}}  \right] \\
        & = (d\omega) \wedge \eta + (-1)^k  \omega \wedge (d \eta)  . 
    \end{align*}        
    \end{solution}   


\begin{proposition}[Nilpotency]
    $d^2 = 0$.
\end{proposition}

\begin{solution}
We first check it for functions. 
\begin{align*}
    d^2 f (X,Y) & = X (df (Y)  )  - Y (df (X)) - df ([X,Y]) \\
    & = X Yf - YXf - [X,Y] f \\
    & = 0 . 
\end{align*}
Then, working in local coordinates, if
\[     \omega = f dx^{i_1} \wedge \cdots \wedge dx ^{i_k}  ,             \]
then
\begin{align*}
    d^2 \omega & = d\left[  df \wedge  dx^{i_1} \wedge \cdots \wedge dx ^{i_k}  \right] \\
    & = ( d^2 f ) \wedge  dx^{i_1} \wedge \cdots \wedge dx ^{i_k}  \\
    & = 0 . 
\end{align*}
\end{solution}


\begin{example}[Classical calculus in $\mathbb{R}^3$]
In $\mathbb{R}^3$, we identify the $1$-form
\[  \eta =   P   dx +  Q  dy + R   dz       \]
with the vector field 
\[  F_{\eta} : =   \langle P, Q, R \rangle    .  \]
In particular, for a function $f \in \Omega ^0 (M) = C^{\infty}(M)$,  
\[   F_{df}  = \nabla f .               \]
We also identify the $2$-form
\[  \omega =  \alpha dy \wedge dz  + \beta dz \wedge dx  + \gamma dx \wedge dy                        \]
with the vector field 
\[   F^{\omega} :  = \langle \alpha  , \beta , \gamma \rangle   .      \]
Recall that the curl of $F_{\eta}$ is given by 
\[    \text{curl} (F_{\eta}) = \langle \tfrac{\partial R}{\partial y} - \tfrac{\partial Q}{\partial z} , \tfrac{\partial P}{\partial z} - \tfrac{\partial R}{\partial x} , \tfrac{\partial Q}{\partial x}  - \tfrac{\partial P}{\partial y} \rangle .          \]
On the other hand, 
\begin{align*}
     d \eta & = \tfrac{\partial P }{\partial y}  \, dy \wedge dx + \tfrac{\partial P }{\partial z} \, dz \wedge dx  + \tfrac{\partial Q }{\partial x} \, dx \wedge dy + \tfrac{\partial Q }{\partial z} \, dz \wedge dy   \\
     &  \,\,\,\,\,\,\,\,\,\,\,\,\,\,\,\,\,\,\,\,\,\,\,\,\,\,\,\,\,     + \tfrac{\partial R }{\partial x} \, dx \wedge dz  + \tfrac{\partial R }{\partial y} \, dy \wedge dz \\
     & = \left(  \tfrac{\partial R }{\partial y} -  \tfrac{\partial Q }{\partial z}  \right) \, dy \wedge dz + \left(  \tfrac{\partial P }{\partial z} -  \tfrac{\partial R }{\partial x}  \right)\, dz \wedge dx + \left(  \tfrac{\partial Q }{\partial x} -  \tfrac{\partial P }{\partial y}  \right)\, dx \wedge dy  .
\end{align*}
Therefore, 
\[    F^{d \eta } = \text{curl}(F_{\eta}).                      \]
Recall that the divergence of $F^{\omega}$ is given by 
\[     \text{div} (F^{\omega}) =  \frac{\partial \alpha}{\partial x} +   \frac{\partial \beta }{\partial y }  +  \frac{\partial \gamma }{\partial z}    .            \]
On the other hand, 
\begin{align*}
    d \omega   = \frac{\partial \alpha}{\partial x } \, & dx \wedge dy \wedge dz + \frac{\partial \beta }{\partial y } \, dy \wedge dz \wedge dx +\frac{\partial \gamma }{\partial z } \, dz \wedge dx \wedge dy \\
    & = \left[ \frac{\partial \alpha}{\partial x } + \frac{\partial \beta }{\partial y } +\frac{\partial \gamma }{\partial z }         \right] \, dx \wedge dy \wedge dz \\
    & = [\text{div} (F^{\omega})] \, dx \wedge dy \wedge dz . 
\end{align*}
\end{example}




\end{document}